% 問題を記述するための環境定義
\def\probs{}
\def\none{none}
\newtcolorbox{probbox}[3]{
	breakable,
	enhanced,
	colback=blue!4!white,
	colframe=blue!55!black,
	coltitle=blue!80!black,
	colbacktitle=blue!15!white,
	boxrule=0.5pt,
	arc=1mm,
	left=1.2mm,
	right=1.2mm,
	top=0.8mm,
	bottom=0.8mm,
	fonttitle=\bfseries,
	title={Q.#1\hspace{0.4em}#2\hspace{0.4em}#3}
}
\newenvironment{prob}[3]%
{\bigskip \global\expandafter\def\expandafter\probs\expandafter{\probs \quad Q.#1}%
 \begin{probbox}{#1}{#2}{#3}}
{\end{probbox}\smallskip}
\newcommand{\probsource}[1]{%
 \tcblower
 {\footnotesize 出典: #1}%
}

% 補題を記述するための環境定義
\definecolor{gray}{rgb}{0.75,0.75,0.75}
\newenvironment{subthm}[1]%
{\def\FrameCommand{\textcolor{gray}{\vrule width 5pt} \hspace{3pt}}
 \MakeFramed {\advance\hsize-\width \FrameRestore}
 \hspace{6pt} \textbf{補題#1} \\
}
{\endMakeFramed}
